\documentclass[conference]{IEEEtran}
\usepackage{cite,algorithmic,balance}
\usepackage{amsmath,amssymb,amsfonts,booktabs,graphicx,textcomp,xcolor}
\usepackage[hidelinks]{hyperref}
\newcounter{algorithm}
% All tables span exactly one IEEE column with shared typography.
\newcommand{\papertablestyle}{\centering\small\setlength{\tabcolsep}{3pt}\renewcommand{\arraystretch}{1.05}}
\newcommand{\E}{\mathbb E}
% Natural two-column flow; retain standard IEEEtran typography.
% IEEE template: https://www.overleaf.com/latex/templates/ieee-conference-template/grfzhhncsfqn
% Preserve IEEEtran margins, columns, fonts and spacing.
% Compile from TPEC_conference using the independent conference_refs.bib.
\title{Informative Power-System Probing via
Sequential Bayesian Experimental Design}
% For two authors, remove the final \and and third author block.
\author{\IEEEauthorblockN{First Author}
\IEEEauthorblockA{\textit{Department / Institution}\\
City, Country\\
email@example.com}
\and
\IEEEauthorblockN{Second Author}
\IEEEauthorblockA{\textit{Department / Institution}\\
City, Country\\
email@example.com}
\and
\IEEEauthorblockN{Third Author}
\IEEEauthorblockA{\textit{Department / Institution}\\
City, Country\\
email@example.com}}
\begin{document}
\raggedbottom
\maketitle
\begin{abstract}
Understanding power-system dynamics requires knowledge of latent parameters
such as inertia and effective frequency response. Active probing provides
a practical means of inferring these parameters by applying controlled
inputs and measuring the resulting frequency response. The information
obtained, however, depends on the probing signals selected. This paper
investigates sequential Bayesian optimal experimental design for acquiring
more information about latent parameters with the same number of probes.
We formulate the selection of probe durations as an expected-information-gain
maximization problem, allowing each decision to use observations from
previous probes while accounting for the evolving system state. We apply
and compare six existing approaches: deep adaptive design,
reinforcement-learning-based sequential design, Step-DAD online refinement,
myopic optimization, fixed design, and random design. The evaluation uses
reduced IEEE9 and IEEE14 models under matched probing constraints.
The primary comparison is a contrastive lower bound on information gain
at a fixed number of design steps; training-seed variability and computational
cost are also assessed.
On IEEE14, DAD achieves $4.7576\pm0.0362$ nats at three design
steps, versus $4.2907\pm0.0368$ for Myopic, with mean online
computation of 1.22 ms versus 10.00 s per sequence.
These results support informative probing through offline-trained designs;
online refinement shows no clear additional benefit in this setting.
On IEEE9 at five design steps, Step-DAD achieves $2.6621\pm0.0145$ nats, compared with DAD at $2.6351\pm0.0179$ nats, a modest mean gain accompanied by additional online computation.
\end{abstract}
\begin{IEEEkeywords}
Bayesian experimental design, power-system identification, adaptive probing,
expected information gain, rate of change of frequency.
\end{IEEEkeywords}

\section{Introduction}
Inertia and effective frequency response govern how a power system reacts
to disturbances. Designed excitation can provide information about both,
motivating probing-based parameter estimation~\cite{peng2024probing}.
Here the question is how to choose successive probe durations when previous
observations inform the next choice and the physical system does not reset.
The reduced dynamics and scalar sensor used below are study assumptions;
they are not a reproduction of a hardware-validated identification system.

Deep adaptive design (DAD) amortizes design optimization into an offline-trained policy and supplies
sequential contrastive information bounds~\cite{foster2021dad}.
Reinforcement-learning-based sequential Bayesian optimal experimental
design (RL-sBOED) learns sequential design policies through simulated rewards
without simulator derivatives~\cite{blau2022rlsboed}. Step-DAD is a semi-amortized method that refines an
offline policy after observations arrive~\cite{hedman2025stepdad}.
We adapt these families to ordered, state-carrying power-grid experiments.
Our contribution is the common application formulation and comparison,
including the information--online-computation trade-off, rather than a claim
that these policy families are new or that adaptivity guarantees improvement.

We compare information gained at equal design-step budgets on IEEE9 and
IEEE14, examining the benefit of adaptive decisions and planning over the
complete sequence. Online runtime is a separate trade-off.

\section{Problem Formulation}\label{sec:problem}
We seek probe durations that reveal the unknown inertia and effective
frequency-response coefficients with a fixed budget of $T$ experiments.
Each experiment supplies one observation, which can inform the next
duration. The model below specifies what is unknown, how a probe produces
data, and which information objective the design policy maximizes.

\subsection{Unknown Parameters and Grid Dynamics}
Let $m$ denote the number of retained machines and define
$\theta=(M_1,\ldots,M_m,K_1,\ldots,K_m)$.
The parameters remain fixed throughout an experimental sequence, with
independent priors $M_i\sim\mathcal U(0.7\bar M_i,1.3\bar M_i)$ and
$K_i\sim\mathcal U(0.1\bar K_i,\bar K_i)$. The reduced dynamics are
\begin{align}
\dot\delta_i &= \omega_i,\label{eq:angle}\\
M_i\dot\omega_i &= v_i-\left(D_i+\frac{K_i}{2\pi}\right)\omega_i
 \nonumber\\[-1pt]
&\quad-\sum_{j=1}^{m}B_{ij}\sin(\delta_i-\delta_j).
\label{eq:swing}
\end{align}
Here $\delta_i$ is rotor angle, $\omega_i$ is angular-speed deviation,
and $\Delta f_i=\omega_i/(2\pi)$ is frequency deviation in Hz.
The input $v_i$ is probing power; $D_i$ is known damping.
The units of $M_i$, $K_i$ and $D_i$ are pu$\,\mathrm{s}^2$/rad,
pu/Hz and pu$\,\mathrm{s}$/rad, respectively, on a 100-MVA base
at nominal frequency 60 Hz. All initial angles and speed deviations
are zero. Net mechanical injections are zero in this deviation model,
so the initial unforced state is an equilibrium.

IEEE9 retains buses 1, 2 and 3; IEEE14 retains buses 1, 2, 3, 6 and 8.
Reactances from the test cases~\cite{zimmerman2011matpower} define a
lossless Laplacian $L$. Eliminating algebraic buses $l$ and retaining
machines $g$ gives $L_{\mathrm r}=L_{gg}-L_{gl}L_{ll}^{-1}L_{lg}$,
with $B_{ij}=-(L_{\mathrm r})_{ij}$ for $i\ne j$ and $B_{ii}=0$.
Resistance, voltage dynamics and transformer tap effects are omitted.
Table~\ref{tab:physical} gives the prior scales and fixed damping,
rounded for display; configurations retain full precision.
Training and evaluation draw independent systems from the same prior.
\begin{table}[t]
\caption{Machine-specific prior scales and fixed damping.}
\label{tab:physical}
\papertablestyle
\begin{tabular*}{\columnwidth}{@{\extracolsep{\fill}}llrrr@{}}
\toprule Network & Bus & $\bar M_i$ & $\bar K_i$ & $D_i$\\\midrule
IEEE9 & 1 & 0.1254141 & 0.2937500 & 0.0125414\\
 & 2 & 0.0339531 & 0.1312500 & 0.0033953\\
 & 3 & 0.0159685 & 0.0750000 & 0.0015969\\
IEEE14 & 1 & 0.0212207 & 0.0784314 & 0.0021221\\
 & 2 & 0.0344836 & 0.1274510 & 0.0034484\\
 & 3, 6, 8 & 0.0265257 & 0.0980392 & 0.0026526\\
\bottomrule
\end{tabular*}
\end{table}

\subsection{Experiments and Conditional Observations}
At design step $t$, the input at bus 1 is a raised-cosine pulse with
duration $\tau_t$ and fixed amplitude $A=0.05$ pu. In local window time
$r\in[0,W]$, it is
\begin{equation}
v_{1,t}(r)=
\begin{cases}
\dfrac A2[1-\cos(2\pi r/\tau_t)],&0\le r\le\tau_t,\\
0,&\tau_t<r\le W.
\end{cases}
\label{eq:probe}
\end{equation}
All other probing inputs are zero. Every window lasts $W=3.5$ s.
Writing $\mathbf x=(\boldsymbol\delta,\boldsymbol\omega)$, the
window-end state is $\mathbf x_t=\Phi_\theta(\mathbf x_{t-1},\tau_t;W)$,
with $\mathbf x_0=\mathbf0$. The next window starts immediately from
$\mathbf x_t$; the system is initialized only once.

Let $f_{1,t}(r;\theta,\tau_{1:t})$ denote the simulated frequency
deviation within window $t$, including the state left by earlier probes.
The sensor measures signed endpoint rate of change of frequency (RoCoF):
\begin{align}
g_t(\theta,\tau_{1:t})
 &=\frac{f_{1,t}(W)-f_{1,t}(W-\Delta r)}{\Delta r},\label{eq:sensor}\\
y_t&=g_t(\theta,\tau_{1:t})+e_t,
\quad e_t\sim\mathcal N(0,\sigma^2).\label{eq:noise}
\end{align}
The frequency arguments in the numerator are suppressed for readability;
$\Delta r=0.025$ s and $\sigma=0.005$ Hz/s. Noise is independent across
windows. Thus $p(y_t\mid\theta,h_{t-1},\tau_t)$ is Gaussian with mean
$g_t$ and variance $\sigma^2$. It depends on the complete duration
history through the carried physical state. Observations arrive at window
end; decision latency is idealized as zero in simulated time.

\subsection{Feasible Adaptive Designs}
The ordered history is $h_t=((\tau_k,y_k))_{k=1}^{t}$, with
$h_0=\varnothing$. A causal policy chooses $\tau_t$ from $h_{t-1}$;
it cannot observe $\theta$ or the true physical state.
Durations must lie in $[0.2,3]$ s and increase by at least
$\varepsilon=0.01$ s. Define
\begin{equation}
a_t=\begin{cases}0.2,&t=1,\\\tau_{t-1}+\varepsilon,&t>1,\end{cases}
\qquad b_t=3-(T-t)\varepsilon.
\label{eq:bounds}
\end{equation}
All methods use $\tau_t=a_t+(b_t-a_t)u_t$, where $u_t\in[0,1]$.
This map reserves room for the remaining design steps without sorting
completed histories or discretizing durations. Let $\Pi$ be the class
of causal policies satisfying these constraints; any policy randomization
is independent of the unknown parameters.

\subsection{Information Objective and Its Approximation}
After observing $h_t$, Bayesian updating gives
\begin{equation}
p(\theta\mid h_t)\propto p_0(\theta)
\prod_{k=1}^{t}p(y_k\mid\theta,h_{k-1},\tau_k).
\label{eq:posterior}
\end{equation}
Each candidate parameter vector propagates its own state under the
observed durations when evaluating this likelihood. Policy action
probabilities cancel in the posterior because, conditional on the
observed history, actions contain no additional access to $\theta$.

The design objective is the expected information gained about $\theta$
after all $T$ observations:
\begin{align}
J_T(\pi)&=\E_{\theta,h_T\mid\pi}
 \left[\log\frac{p(\theta\mid h_T)}{p_0(\theta)}\right],
 \label{eq:eig}\\
\pi^\star&\in\operatorname*{arg\,max}_{\pi\in\Pi}\,J_T(\pi).
\label{eq:design_problem}
\end{align}
The expectation averages parameter, observation and policy randomness.
This is an objective for the complete sequence: every method receives
the same $T$ probes, while its offline and online computational costs
are assessed separately.

The marginal likelihood required by exact EIG is not tractable here.
We therefore use the sequential prior-contrastive (sPCE) lower
bound~\cite{foster2021dad}. Draw $\theta^{(0)}$ from the prior to
generate a trajectory and draw $L$ independent prior contrasts
$\theta^{(1)},\ldots,\theta^{(L)}$. Define their trajectory likelihoods
and the resulting score by
\begin{align}
q_t^{(n)}&=\prod_{k=1}^{t}
 p(y_k\mid\theta^{(n)},h_{k-1},\tau_k),\\
\widehat I_t&=\log\frac{q_t^{(0)}}
 {\frac{1}{L+1}\sum_{n=0}^{L}q_t^{(n)}}.
\label{eq:spce}
\end{align}
Likelihoods are computed in the log domain. The mean score estimates
a lower bound on $J_T$, whereas an individual score can be negative.
The score cannot exceed $\log(L+1)$: with the conference setting
$L=128$, this ceiling is 4.8598 nats. Near-ceiling scores can conceal
differences in true EIG. The generating parameter appears in the
evaluator's denominator only; it is never supplied to an online planner.

\section{Compared Design Methods}
All six methods select the same continuous probe duration from the
feasible interval in~\eqref{eq:bounds}. They differ in whether they use
observations, optimize online, or learn a policy offline. Algorithm~\ref{alg:probing}
summarizes their evaluation procedure; the following subsections specify
how each decision is obtained. These are application-specific adaptations
of existing methods, not a new design algorithm.

\begin{figure}[t]
\small
\refstepcounter{algorithm}\label{alg:probing}
\hrule\smallskip
\textbf{Algorithm \thealgorithm: Evaluate one probe sequence}
\smallskip\hrule\smallskip
\begin{algorithmic}[1]
\REQUIRE Method $\mathsf M$, prior $p_0$, steps $T$, selected offline model if needed
\STATE Set $h_0=\varnothing$; initialize the physical system once.
\IF{$\mathsf M\in\{\mathrm{Myopic},\mathrm{Step\mbox{-}DAD}\}$}
\STATE Initialize 128 prior particles, their states and equal weights.
\ENDIF
\IF{$\mathsf M=\mathrm{Step\mbox{-}DAD}$}
\STATE Set episode policy $\phi$ to the matching DAD checkpoint.
\ENDIF
\FOR{$t=1,\ldots,T$}
\IF{$t>1$ and $\mathsf M\in\{\mathrm{Myopic},\mathrm{Step\mbox{-}DAD}\}$}
\STATE Propagate particles under $\tau_{t-1}$; update their weights with $y_{t-1}$.
\ENDIF
\IF{$\mathsf M=\mathrm{Random}$}
\STATE Draw $u_t\sim\mathcal U(0,1)$ and map it to $\tau_t\in[a_t,b_t]$.
\ELSIF{$\mathsf M=\mathrm{Myopic}$}
\STATE Search feasible durations for maximum expected one-step information.
\ELSE
\IF{$\mathsf M=\mathrm{Step\mbox{-}DAD}$ and $t>1$}
\STATE Refine $\phi$ on posterior future simulations for four updates; keep the conditional-validation-best candidate, including the unchanged policy.
\ENDIF
\STATE Evaluate the Fixed/DAD/RL-sBOED model or $\phi$; map its output to $\tau_t\in[a_t,b_t]$.
\ENDIF
\STATE Apply the pulse; carry the physical state to window end and observe $y_t$.
\STATE Append $(\tau_t,y_t)$ to form $h_t$.
\ENDFOR
\ENSURE History $h_T$; score terminal sPCE separately.
\end{algorithmic}
\smallskip\hrule
\end{figure}

\subsection{Nonadaptive Designs: Random and Fixed}
Random samples a unit action uniformly and applies the common feasible
duration map. This is uniform sampling within the remaining interval,
not over complete ordered sequences. It requires no offline training
and ignores observations when selecting durations.

Fixed optimizes an entire ordered sequence offline using terminal sPCE.
It compares short, central, long and spread interior initializations
using validation, then optimizes the selected initialization. Each
training seed yields a separately fitted sequence. During evaluation,
Fixed executes that sequence without adapting to observations.

\subsection{Online One-Step Design: Myopic}
Myopic maintains a weighted particle approximation to the parameter
posterior. Before each decision after the first, particles propagate
their own states under the preceding duration and update likelihood
weights using the latest observation. Continuous search then maximizes
expected information from the next probe only, using 16 proposals,
three search rounds and 32 simulated observations per proposal.
Candidate durations obey the same constraints as all other methods.
No offline policy is trained.

\subsection{Offline-Trained Adaptive Policies: DAD and RL-sBOED}
Both policies receive the stage index and ordered slots containing past
durations, scaled observations and availability masks; future slots are
masked. They do not receive true parameters or physical states.
At evaluation, a forward pass produces the action without online
optimization or an explicit particle-posterior update.

DAD uses a deterministic two-layer, 64-unit tanh network, a trainable
stage bias and a sigmoid action map. Offline training maximizes terminal
sPCE by differentiating through the complete simulated sequence,
including observations and carried states. In the joint Fixed/DAD
campaign, validation selects between random and interior Fixed-based
initialization. The reported DAD stage timer excludes the preceding
Fixed training; the full workflow incurs both costs.

RL-sBOED uses a REDQ-style actor--critic adaptation with ten critics
and five replay updates per trajectory batch. The actor has two
128-unit ReLU layers. Its rewards are
$r_t=\widehat I_t-\widehat I_{t-1}$, with $\widehat I_0=0$ and
discount $\gamma=1$, so their sum is terminal sPCE. Training therefore
accounts for future design steps without differentiating through the
simulator. Entropy regularization uses normalized tanh actions and
target entropy $-1$; reported information scores exclude entropy bonuses.
Evaluation uses the actor's deterministic output.

\subsection{Online Policy Refinement: Step-DAD}
Step-DAD begins each episode with the matching selected DAD checkpoint
and uses it unchanged for the first probe. Before each subsequent
decision, it assimilates the latest observation into its 128-particle
posterior. It then samples candidate parameters and their carried
states from that posterior to simulate the remaining design steps.
These future simulations condition on the actual observed history.

At each refinement, a fresh Adam optimizer performs four gradient
updates with batch size 32 and learning rate $10^{-3}$. A fixed
conditional validation batch compares the unchanged policy and all
updated candidates; the best is retained. Retained weights carry to
the next design step within the episode, but not to another episode.
The total refinement budget is $4(T-1)$ updates: 8, 12 and 16 for
$T=3,4,5$. The planner never receives the evaluator's true parameters.

For both Myopic and Step-DAD, assimilation is performed only when a
subsequent decision requires it. No final planner update is needed
after $y_T$ for this EIG experiment; terminal information is computed
separately by the evaluator using~\eqref{eq:spce}.

\section{Experimental Settings}
All six methods are compared at $T=3,4,5$ design steps on IEEE9 and
IEEE14. Each step applies one probe and collects one observation.
Equal design-step counts define the common probe budget. The systems
share sensor and design constraints but use separate physical parameters
and trained policies.

The swing equations are integrated by RK4 with step $1/640$ s.
Numerical state/duration Jacobians use perturbations of $10^{-5}$
and include carried-state derivatives. The verification suite compares
policy and conditional-refinement gradients with end-to-end finite
differences.

Offline training uses 2,000 updates with batch size 32 and three independent
training seeds (101, 202 and 303). The DAD and Fixed learning rate is
$10^{-3}$. Validation uses 128 systems every 100 updates. Training,
validation and evaluation use 128 contrastive samples for the sPCE
objective. Myopic and Step-DAD use 128 particles for online posterior
planning; Step-DAD performs four refinement updates after each observation
when design steps remain. Evaluation uses seeds 1001, 1002 and 1003,
with 512 simulated experimental sequences per evaluation seed.
Validation selects checkpoints without gradient updates on validation cases.
Each trained policy is evaluated on 1,536 episodes. Parameter and
measurement-noise draws are paired across methods; the online planner uses
a separate random stream. Random/Myopic have no training-seed dependence
under these streams, so identical evaluations can be shared across training
seeds. Fixed remains independently optimized for each training seed.

For trained methods, we average evaluation-seed means within each training
seed and report mean $\pm$ sample standard deviation across three trainings.
Random/Myopic report mean $\pm$ sample SD of evaluation-seed means,
not episode-level dispersion. The conference evaluation protocol uses
three seeds (1001--1003); all required runs are complete.
Evaluation seed 1001 informed development and is not untouched test data.

The project uses NVIDIA GeForce RTX 5080 and RTX 4090 workstations
and NVIDIA A100 cluster GPUs for development and experiments.
The reported grid timing runs use RTX 5080 and A100 GPUs; the RTX 4090
is the development platform. Computation is reported in
Section~\ref{sec:cost}. IEEE9 offline times average runs across RTX 5080
and A100 GPUs, while IEEE14 uses A100s. Online timings use RTX 5080 for IEEE9, except the migrated Step-DAD
$T=5$ seed 303 on A100; its reported mean combines those runs.
IEEE14 online timings use A100. These are recorded workflow
costs, not a controlled comparison of runtime between networks.
Configurations, source hashes and checkpoints preserve provenance.
Code: \url{https://github.com/alexschmidt123/objective_power_grid}.

% Complete results under the accepted three-evaluation-seed grid protocol.
\section{Results and Discussion}
\label{sec:results}

\subsection{Information Versus Number of Design Steps}
\begin{table}[!ht]\caption{IEEE9 terminal sPCE (nats), mean $\pm$ sample SD. Trained methods: three training-seed means; Random/Myopic: three evaluation-seed means. $L=128$; ceiling 4.8598.}\label{tab:ieee9}\papertablestyle
\begin{tabular*}{\columnwidth}{@{\extracolsep{\fill}}lccc@{}}\toprule Method & $T=3$ & $T=4$ & $T=5$\\\midrule
Random & $0.4321\pm0.0310$ & $0.4924\pm0.0342$ & $0.5392\pm0.0471$\\
Fixed & $2.0112\pm0.0116$ & $2.3192\pm0.0004$ & $2.5066\pm0.0013$\\
Myopic & $1.9736\pm0.0438$ & $2.1707\pm0.0366$ & $2.2407\pm0.0547$\\
DAD & $2.1366\pm0.0161$ & $2.4095\pm0.0080$ & $2.6351\pm0.0179$\\
RL-sBOED & $2.1257\pm0.0027$ & $2.3963\pm0.0333$ & $2.5307\pm0.0114$\\
Step-DAD & $2.1471\pm0.0110$ & $2.4382\pm0.0048$ & $2.6621\pm0.0145$\\
\bottomrule\end{tabular*}\end{table}

\begin{table}[!ht]
\caption{IEEE14 terminal sPCE (nats), mean $\pm$ sample SD.
Trained methods use three training-seed means; Random/Myopic use
three evaluation-seed means. $L=128$; ceiling 4.8598.}
\label{tab:ieee14}
\papertablestyle
\begin{tabular*}{\columnwidth}{@{\extracolsep{\fill}}lccc@{}}
\toprule Method & $T=3$ & $T=4$ & $T=5$\\\midrule
Random & $0.7541\pm0.0047$ & $0.8360\pm0.0297$ & $0.8464\pm0.0252$\\
Fixed & $4.2623\pm0.0402$ & $4.4946\pm0.0345$ & $4.6624\pm0.0083$\\
Myopic & $4.2907\pm0.0368$ & $4.3990\pm0.0177$ & $4.4641\pm0.0280$\\
DAD & $4.7576\pm0.0362$ & $4.8325\pm0.0047$ & $4.8443\pm0.0026$\\
RL-sBOED & $4.7445\pm0.0022$ & $4.8275\pm0.0112$ & $4.8383\pm0.0108$\\
Step-DAD & $4.7527\pm0.0275$ & $4.8297\pm0.0067$ & $4.8446\pm0.0042$\\
\bottomrule\end{tabular*}\end{table}
On IEEE9, DAD exceeds Fixed by 0.1253, 0.0903 and 0.1286 nats
at $T=3,4,5$, respectively. The observed Step-DAD gains over DAD
are 0.0105, 0.0287 and 0.0269 nats at $T=3,4,5$, respectively.
These modest mean gains come with substantially greater online cost. IEEE9 scores stay farther from
the estimator ceiling than IEEE14 scores. Runtime is reported within
each system in Tables~\ref{tab:ieee9_online} and~\ref{tab:ieee14_online}.



For IEEE14, Table~\ref{tab:ieee14} shows higher mean sPCE for
DAD, RL-sBOED and Step-DAD than for Random, Fixed or Myopic at
every design-step count. At $T=3$, DAD exceeds Myopic by 0.4669 nats
and Fixed by 0.4953 nats. The comparison fixes the number of probes,
so the gain comes from selecting more informative experiments within
the same probe budget. Fixed exceeds Myopic at $T=4,5$; observation
adaptation alone therefore does not guarantee a stronger design.

The DAD-family means are close to one another. Step-DAD changes the
DAD mean by $-0.0049$, $-0.0028$ and $+0.0003$ nats for
$T=3,4,5$, respectively. These small differences do not establish an
advantage for online refinement. At $T=5$, DAD lies only 0.0155 nats
below the 4.8598-nat estimator ceiling. Finite-contrast saturation can
compress method differences; these scores neither measure exact EIG
nor establish a reliable ranking within the DAD family. The SDs
describe distinct sources of variation: retraining for learned/fitted
methods and repeated evaluation for Random/Myopic.
\subsection{Measured Computational Cost}\label{sec:cost}
Online cost is the wall-clock time summed over all $T$ decisions in
one experimental sequence, averaged over evaluation episodes and then
training seeds where applicable. It includes action selection, posterior
updates needed for subsequent decisions, Myopic search and Step-DAD
refinement. It excludes the simulated physical response, observation
acquisition and evaluator sPCE scoring. These are computational timings,
not end-to-end physical experiment durations.





\begin{table}[!ht]
\caption{IEEE9 mean offline training time (hours per training seed).}
\label{tab:ieee9_offline}
\papertablestyle
\begin{tabular*}{\columnwidth}{@{\extracolsep{\fill}}lccc@{}}
\toprule Method & $T=3$ & $T=4$ & $T=5$\\\midrule
Random & 0.000 & 0.000 & 0.000\\
Fixed & 0.960 & 1.396 & 1.827\\
Myopic & 0.000 & 0.000 & 0.000\\
DAD & 0.957 & 1.393 & 1.831\\
RL-sBOED & 0.120 & 0.152 & 0.180\\
Step-DAD & \multicolumn{3}{c}{Reuses matching DAD model}\\
\bottomrule\end{tabular*}
\end{table}

\begin{table}[!ht]
\caption{IEEE14 mean offline training time (hours per training seed).}
\label{tab:ieee14_offline}
\papertablestyle
\begin{tabular*}{\columnwidth}{@{\extracolsep{\fill}}lccc@{}}
\toprule Method & $T=3$ & $T=4$ & $T=5$\\\midrule
Random & 0.000 & 0.000 & 0.000\\
Fixed & 1.877 & 2.750 & 3.691\\
Myopic & 0.000 & 0.000 & 0.000\\
DAD & 1.894 & 2.725 & 3.626\\
RL-sBOED & 0.170 & 0.204 & 0.248\\
Step-DAD & \multicolumn{3}{c}{Reuses matching DAD model}\\
\bottomrule\end{tabular*}
\end{table}

\begin{table}[!ht]
\caption{IEEE9 mean online computation (seconds per complete sequence).}
\label{tab:ieee9_online}
\papertablestyle
\begin{tabular*}{\columnwidth}{@{\extracolsep{\fill}}lccc@{}}
\toprule Method & $T=3$ & $T=4$ & $T=5$\\\midrule
Random & 0.00007 & 0.00009 & 0.00012\\
Fixed & 0.00052 & 0.00062 & 0.00077\\
Myopic & 9.30817 & 12.57734 & 15.89200\\
DAD & 0.00054 & 0.00068 & 0.00084\\
RL-sBOED & 0.00049 & 0.00063 & 0.00076\\
Step-DAD & 6.34247 & 15.80823 & 24.34984\\
\bottomrule\end{tabular*}
\end{table}

\begin{table}[!ht]
\caption{IEEE14 mean online computation (seconds per complete sequence).}
\label{tab:ieee14_online}
\papertablestyle
\begin{tabular*}{\columnwidth}{@{\extracolsep{\fill}}lccc@{}}
\toprule Method & $T=3$ & $T=4$ & $T=5$\\\midrule
Random & 0.00016 & 0.00021 & 0.00031\\
Fixed & 0.00116 & 0.00148 & 0.00275\\
Myopic & 9.99968 & 13.43718 & 16.93535\\
DAD & 0.00122 & 0.00142 & 0.00193\\
RL-sBOED & 0.00168 & 0.00331 & 0.00271\\
Step-DAD & 9.11275 & 23.24278 & 44.69261\\
\bottomrule\end{tabular*}
\end{table}



Tables~\ref{tab:ieee9_offline}--\ref{tab:ieee14_online} separate offline training
from online computation across all three design-step counts.
The offline timer includes initialization selection, optimization,
periodic validation and checkpoint writes. DAD's reported stage time
excludes the separately incurred Fixed training used for initialization;
a joint Fixed--DAD workflow must include both costs. Step-DAD reuses
that DAD checkpoint and incurs its additional work online.

At $T=5$, DAD requires 1.93 ms online per sequence, compared with
16.94 s for Myopic and 44.69 s for Step-DAD
(Table~\ref{tab:ieee14_online}). Thus the near-identical DAD/Step-DAD scores
have markedly different computational costs. Offline training amortizes
design computation over repeated experiments, but this benefit depends
on how often a trained policy is reused and on the physical deadline.
The current simulation assumes zero decision latency; it does not
demonstrate real-time feasibility for the online planners.

\subsection{Optimization and Estimator Diagnostics}


Validation histories are retained in the companion result collection.
Evaluation uses each policy's best validation checkpoint. DAD can
initialize from the separately trained Fixed design; its update axis
omits that preceding cost, reported separately in Tables~\ref{tab:ieee9_offline} and~\ref{tab:ieee14_offline}.

For IEEE14 training seed 101, the DAD and RL-sBOED checkpoints
occur at updates 1,200 and 1,000, respectively, whereas Fixed selects
update 2,000. Later updates need not improve validation utility.
Neither a fixed update budget nor the recorded validation histories establish convergence.
Contrast-count and planner-particle sensitivity remain to be tested;
being below the sPCE ceiling does not establish estimator accuracy.

\subsection{Scope and Limitations}
The reduced lossless network, independent priors and scalar sensor restrict
the conclusions to the stated simulation model. Parameter distributions
match between training and evaluation; robustness to prior shift, sensor
bias and unmodeled dynamics is not established by this protocol. Identical
optimization budgets do not imply equal convergence or equal computational
effort. In particular, a fixed 2,000-update budget needs interpretation
alongside its validation history.

The reported score is a finite-contrast lower bound, so estimator saturation
can conceal information differences. Online inference also uses a finite
particle population. These approximation errors differ from ODE integration
error and should not be conflated. Information gain alone establishes
neither terminal-control safety nor field performance.

\subsection{Supplementary SIR-ODE Benchmark}
SIR-ODE provides an additional benchmark for sequential measurement
selection. Its susceptible--infected--recovered model uses a population
of 500 with two initially infected individuals. Designs select increasing
measurement times from 100 candidates; observations are noisy infected
counts with Gaussian standard deviation 1.0. This historical experiment
uses a precomputed simulation bank and 1,024 posterior-support particles.
It reports posterior entropy reduction, with prior entropy
$\log(1024)=6.9315$ nats, rather than the grid sPCE estimator.

The archive contains three training seeds and five evaluation seeds,
with 512 systems per evaluation. Random uses 32 design replicates per
system, averaged within system before aggregation. Historical policy
training uses a behavioral-cloning warm start and different optimizers
from the grid experiment; checkpoint selection also includes a configured
sequence-diversity preference. SIR is therefore a supplementary historical
comparison, not a matched replication of the grid implementations.
\begin{table}[!ht]\caption{SIR-ODE entropy reduction (nats), mean $\pm$ sample SD. DAD/RL-sBOED/Step-DAD: three training-seed means. Random/Myopic/Fixed: five evaluation-seed means after deduplicating identical copies across training seeds.}\label{tab:sir}\papertablestyle
\begin{tabular*}{\columnwidth}{@{\extracolsep{\fill}}lccc@{}}\toprule Method & $T=3$ & $T=4$ & $T=5$\\\midrule
Random & $3.7783\pm0.0167$ & $3.8593\pm0.0167$ & $3.9241\pm0.0167$\\
Fixed & $6.2816\pm0.0101$ & $6.4501\pm0.0091$ & $6.5360\pm0.0044$\\
Myopic & $6.7108\pm0.0125$ & $6.7437\pm0.0124$ & $6.7596\pm0.0135$\\
DAD & $6.8394\pm0.0087$ & $6.8591\pm0.0039$ & $6.8658\pm0.0077$\\
RL-sBOED & $6.8494\pm0.0029$ & $6.8674\pm0.0020$ & $6.8814\pm0.0015$\\
Step-DAD & $6.8438\pm0.0095$ & $6.8679\pm0.0011$ & $6.8761\pm0.0109$\\
\bottomrule\end{tabular*}\end{table}

The DAD-family means exceed the baselines for all three design-step
counts, but lie close to the finite-particle entropy ceiling. Small
differences within that family do not establish superiority. The SIR
Fixed sequence is calibrated once, so its uncertainty is across evaluation
seeds; grid Fixed is independently fitted for each training seed.
SIR timing includes observation lookup and posterior updates, with
incompletely recorded historical hardware. It is not pooled with the
grid decision timings. Full configurations, model files and per-run
timings are retained in the companion result collection.

\balance
\section{Conclusion}
This paper formulates sequential power-grid probing as an information-driven
design problem in which each observation informs the next duration and the
physical state carries between probes. The six-method comparison shares a
continuous feasible action map, endpoint-frequency likelihood and paired
evaluation protocol. Separating training-seed variation, finite-contrast
information and online decision time makes the costs and benefits of
adaptation explicit. The resulting framework supports reproducible
assessment on reduced IEEE9 and IEEE14 systems without equating a high
contrastive score with exact information gain or operational readiness.
On IEEE14, the offline-trained DAD family yields higher mean sPCE
than the three baselines at equal design-step counts. Four-update
Step-DAD provides no clear gain over DAD while adding substantial
online computation. On IEEE9, Step-DAD provides modest mean gains
over DAD at all three design-step counts. Near-ceiling IEEE14 scores
and development use of evaluation seed 1001 limit the conclusions.
Extensions include model mismatch and decision latency.
Mixture-of-experts sequential design (MoE-sBOED) is optional future work
outside the present six-method comparison.

\bibliographystyle{IEEEtran}
\bibliography{conference_refs}
\end{document}
